\documentclass[12pt,addpoints]{exam}
\usepackage[margin=1in]{geometry}
\usepackage{amsmath,amssymb,amsthm}
\usepackage{graphicx}
\usepackage{fancyhdr}
\usepackage{enumitem}
\usepackage{multicol}

\pagestyle{fancy}
\lhead{\textbf{Geometry Honors --- Midterm Exam}}
\rhead{Name: \underline{\hspace{3cm}} \quad Date: \underline{\hspace{3cm}}}
\cfoot{Page \thepage\ of \numpages}

\setlength{\headheight}{15pt}
\setlength{\parindent}{0pt}

\pointsinrightmargin
\runningheadrule

\newcommand{\answerbox}[1]{\underline{\hspace{#1}}}
\newcommand{\deg}{{}^\circ}

\begin{document}

% ============================================================
% TITLE PAGE
% ============================================================
\begin{center}
\vspace*{1.5in}
{\Huge\bfseries Geometry Honors --- Midterm Examination\par}
\vspace{0.3in}

{\large\bfseries Topics:} Logical Reasoning $\cdot$ Proofs $\cdot$ Parallel Lines $\cdot$ Triangles $\cdot$ Quadrilaterals $\cdot$ Transformations $\cdot$ Similarity $\cdot$ Right Triangles \& Trigonometry $\cdot$ Circles $\cdot$ Area \& Volume

\vspace{0.2in}
{\large\bfseries Total Points:} 150 \qquad {\bfseries Time Allowed:} 180 minutes \qquad {\bfseries Show all work.}

\vspace{0.5in}
Name: \underline{\hspace{4cm}} \qquad Date: \underline{\hspace{3cm}} \qquad Period: \underline{\hspace{2cm}}
\end{center}

\vspace{0.4in}

\begin{center}
\framebox[\linewidth]{%
\parbox{0.95\linewidth}{%
\textbf{Instructions:} This exam has \textbf{5 sections} covering the Geometry Honors curriculum. Show all work for full credit. Justify answers using definitions, postulates, and theorems. A calculator is permitted for Sections IV--V. No electronic devices for Sections I--III.
}}
\end{center}

\vspace{0.3in}
\begin{center}
\begin{tabular}{|l|c|c|}
\hline
\textbf{Section} & \textbf{Problems} & \textbf{Points} \\ \hline
I. Multiple Choice & 1--15 & 30 \\ \hline
II. Short Answer & 16--25 & 30 \\ \hline
III. Proofs \& Constructions & 26--30 & 30 \\ \hline
IV. Applications \& Word Problems & 31--38 & 30 \\ \hline
V. Challenge Problems & 39--40 & 30 \\ \hline
\textbf{Total} & \textbf{40} & \textbf{150} \\ \hline
\end{tabular}
\end{center}

\newpage

% ============================================================
% SECTION I: MULTIPLE CHOICE (15 questions, 2 pts each)
% ============================================================
\section*{Section I: Multiple Choice \hfill (30 points)}
\noindent\textit{Circle the best answer. Each question is worth 2 points.}

\begin{questions}

\question Which of the following is the \textbf{converse} of the statement ``If two lines are parallel, then corresponding angles are congruent''?
\begin{choices}
  \choice If two lines are not parallel, then corresponding angles are not congruent.
  \choice If corresponding angles are congruent, then two lines are parallel.
  \choice If corresponding angles are not congruent, then two lines are not parallel.
  \choice If two lines are parallel, then alternate interior angles are congruent.
\end{choices}

\question In $\triangle ABC$, $m\angle A = 52\deg$ and $m\angle B = 68\deg$. What is $m\angle C$?
\begin{choices}
  \choice $40\deg$
  \choice $52\deg$
  \choice $60\deg$
  \choice $120\deg$
\end{choices}

\question Which set of side lengths does \textbf{not} form a triangle?
\begin{choices}
  \choice $3,\ 4,\ 5$
  \choice $7,\ 10,\ 12$
  \choice $1,\ 2,\ 3$
  \choice $8,\ 8,\ 15$
\end{choices}

\question The contrapositive of ``If a quadrilateral is a square, then it is a rhombus'' is:
\begin{choices}
  \choice If a quadrilateral is a rhombus, then it is a square.
  \choice If a quadrilateral is not a square, then it is not a rhombus.
  \choice If a quadrilateral is not a rhombus, then it is not a square.
  \choice If a quadrilateral is a rhombus, then it is not a square.
\end{choices}

\question Which transformation \textbf{does not} always preserve congruence?
\begin{choices}
  \choice Translation
  \choice Reflection
  \choice Rotation
  \choice Dilation with scale factor $k \neq 1$
\end{choices}

\question In parallelogram $ABCD$, $m\angle A = (3x + 15)\deg$ and $m\angle B = (2x + 5)\deg$. Find $x$.
\begin{choices}
  \choice $16$
  \choice $28$
  \choice $35$
  \choice $48$
\end{choices}

\question Two similar triangles have areas $36~\text{cm}^2$ and $100~\text{cm}^2$. If the shorter base of the smaller triangle is $6~\text{cm}$, what is the corresponding base of the larger triangle?
\begin{choices}
  \choice $10~\text{cm}$
  \choice $15~\text{cm}$
  \choice $18~\text{cm}$
  \choice $25~\text{cm}$
\end{choices}

\question In right $\triangle ABC$ with $\angle C = 90\deg$, $AC = 9$, $BC = 12$. What is $\sin A$?
\begin{choices}
  \choice $\dfrac{3}{5}$
  \choice $\dfrac{4}{5}$
  \choice $\dfrac{3}{4}$
  \choice $\dfrac{5}{4}$
\end{choices}

\question A central angle in a circle measures $120\deg$. What is the measure of the \textbf{intercepted arc}?
\begin{choices}
  \choice $30\deg$
  \choice $60\deg$
  \choice $120\deg$
  \choice $240\deg$
\end{choices}

\question The circumference of a circle is $16\pi$. What is its area?
\begin{choices}
  \choice $8\pi$
  \choice $16\pi$
  \choice $32\pi$
  \choice $64\pi$
\end{choices}

\question Which point of concurrency is the intersection of the three \textbf{perpendicular bisectors} of a triangle?
\begin{choices}
  \choice Incenter
  \choice Centroid
  \choice Circumcenter
  \choice Orthocenter
\end{choices}

\question The volume of a sphere with radius $6~\text{cm}$ is:
\begin{choices}
  \choice $36\pi~\text{cm}^3$
  \choice $72\pi~\text{cm}^3$
  \choice $144\pi~\text{cm}^3$
  \choice $288\pi~\text{cm}^3$
\end{choices}

\question Lines $l \parallel m$ and transversal $t$ intersects them. If $\angle 1 = (5x - 10)\deg$ and $\angle 2 = (3x + 20)\deg$ are alternate exterior angles, find $x$.
\begin{choices}
  \choice $10$
  \choice $15$
  \choice $17.5$
  \choice $20$
\end{choices}

\question What is the equation of a circle centered at $(-3, 4)$ with radius $7$?
\begin{choices}
  \choice $(x + 3)^2 + (y - 4)^2 = 49$
  \choice $(x - 3)^2 + (y + 4)^2 = 49$
  \choice $(x + 3)^2 + (y - 4)^2 = 7$
  \choice $(x - 3)^2 + (y + 4)^2 = 7$
\end{choices}

\question If $\triangle ABC \sim \triangle DEF$ and the ratio of corresponding sides is $3:5$, what is the ratio of their perimeters?
\begin{choices}
  \choice $3:5$
  \choice $9:25$
  \choice $5:3$
  \choice $25:9$
\end{choices}

\end{questions}

\newpage

% ============================================================
% SECTION II: SHORT ANSWER
% ============================================================
\section*{Section II: Short Answer \hfill (30 points)}
\noindent\textit{Show your work. Each question is worth 3 points.}

\begin{questions}

\question Simplify: $\cos^2 \theta + \sin^2 \theta =$ \answerbox{3cm}

\question In the diagram, $\overline{AB}$ is tangent to circle $O$ at point $B$, and radius $OB = 8~\text{cm}$. If $OA = 10~\text{cm}$, find the length of $\overline{AB}$.

\vspace{1.5cm}

\question A regular hexagon has side length $6$. Find its apothem (exact form).

\vspace{1.5cm}

\question The coordinates of $\triangle ABC$ are $A(1, 2)$, $B(4, 2)$, $C(1, 6)$. After reflecting over the $y$-axis, then translating $3$ units right, what are the coordinates of $A''$?

\vspace{1.5cm}

\question In $\triangle PQR$, $S$ and $T$ are midpoints of $\overline{PQ}$ and $\overline{PR}$ respectively. If $QR = 14$, find $ST$.

\vspace{1.5cm}

\question Find the measure of an inscribed angle that intercepts an arc of $230\deg$.

\vspace{1.5cm}

\question A right triangle has legs of length $5$ and $12$. An altitude is drawn to the hypotenuse. Find the length of this altitude (exact or decimal to 2 places).

\vspace{1.5cm}

\question The surface area of a right cylinder with radius $4~\text{cm}$ and height $10~\text{cm}$ is \answerbox{4cm} $\text{cm}^2$.

\question Two secants are drawn from exterior point $P$ to a circle, intercepting arcs of $110\deg$ and $50\deg$. Find $m\angle P$.

\vspace{1.5cm}

\question A regular pyramid has a square base with side $8$ and slant height $7$. Find its lateral surface area.

\vspace{1.5cm}

\end{questions}

\newpage

% ============================================================
% SECTION III: PROOFS \& CONSTRUCTIONS
% ============================================================
\section*{Section III: Proofs \& Constructions \hfill (30 points)}
\noindent\textit{Each question is worth 6 points. Show all reasoning.}

\begin{questions}

\question \textbf{Two-column proof.} Given: $\overline{AB} \cong \overline{CD}$ and $\overline{BC} \cong \overline{DA}$.\\
Prove: Quadrilateral $ABCD$ is a parallelogram.

\vspace{3cm}

\question \textbf{Paragraph proof.} Given: Point $D$ is the midpoint of $\overline{BC}$ in $\triangle ABC$. Line $AD$ is extended to $E$ such that $\overline{AD} \cong \overline{DE}$. Segments $BE$ and $CE$ are drawn.\\
Prove: $BE + CE > BC$.

\vspace{3cm}

\question \textbf{Indirect proof.} Prove: In a right triangle, each acute angle measures less than $90\deg$.

\vspace{3cm}

\question \textbf{Construction.} Using only a compass and straightedge, describe the step-by-step construction of the perpendicular bisector of segment $\overline{AB}$. Then prove that your construction is correct.

\vspace{3cm}

\question \textbf{Coordinate geometry proof.} Given: Points $A(-2, 3)$, $B(4, 7)$, $C(8, 1)$, $D(2, -3)$.\\
Prove that $ABCD$ is a parallelogram by showing both pairs of opposite sides are parallel. Then determine whether it is a rhombus, rectangle, or square, and justify your answer.

\vspace{3cm}

\end{questions}

\newpage

% ============================================================
% SECTION IV: APPLICATIONS \& WORD PROBLEMS
% ============================================================
\section*{Section IV: Applications \& Word Problems \hfill (30 points)}
\noindent\textit{Show all work. Point values are indicated in brackets.}

\begin{questions}

\question[4] A ladder $15~\text{ft}$ long leans against a building. If the foot of the ladder is $9~\text{ft}$ from the base of the building, how high up the wall does the ladder reach?

\vspace{1.5cm}

\question[5] A tree casts a shadow $24~\text{m}$ long. At the same time, a $2~\text{m}$ pole casts a shadow $3~\text{m}$ long. How tall is the tree?

\vspace{1.5cm}

\question[5] A circular fountain has diameter $12~\text{ft}$. A $3~\text{ft}$ wide walkway surrounds it. Find the area of the walkway in terms of $\pi$.

\vspace{1.5cm}

\question[5] A satellite dish is parabolic (modeled as a circular arc). If the dish is $8~\text{ft}$ across and $2~\text{ft}$ deep, find the radius of the circle that models the dish.

\vspace{1.5cm}

\question[5] The centroid of $\triangle ABC$ has coordinates $(4, 5)$. If $A = (1, 2)$ and $B = (7, 8)$, find the coordinates of $C$.

\vspace{1.5cm}

\question[5] A regular octagon is inscribed in a circle of radius $10$. Find the area of the octagon to the nearest whole number.

\vspace{1.5cm}

\question[5] The volume of a cone is $120\pi~\text{in}^3$ and the height is $12~\text{in}$. Find the slant height.

\vspace{1.5cm}

\question[4] In $\triangle ABC$, $\angle C = 90\deg$. The altitude from $C$ to $\overline{AB}$ divides $\overline{AB}$ into segments of length $3$ and $12$. Find the length of this altitude.

\vspace{1.5cm}

\end{questions}

\newpage

% ============================================================
% SECTION V: CHALLENGE PROBLEMS
% ============================================================
\section*{Section V: Challenge Problems \hfill (30 points)}
\noindent\textit{These problems require synthesis of multiple concepts. Show all reasoning. Each question is worth 15 points.}

\bigskip

\noindent\textbf{Problem 39.} \textbf{Geometry of a regular polygon and trigonometry.}

A regular dodecagon (12-sided polygon) is inscribed in a circle of radius $R$.

\begin{enumerate}[label=(\alph*),leftmargin=1.5cm]
  \item Derive a formula for the side length of the dodecagon in terms of $R$.
  \item Derive a formula for the area of the dodecagon in terms of $R$.
  \item If $R = 10$, find the area to the nearest tenth.
  \item Explain why the area of a regular $n$-gon inscribed in a circle of radius $R$ approaches $\pi R^2$ as $n \to \infty$.
\end{enumerate}

\vspace{3cm}

\noindent\textbf{Problem 40.} \textbf{Comprehensive proof and application.}

In $\triangle ABC$, $AB = AC$. Let $D$ be a point on $\overline{BC}$. The circumcircle of $\triangle ABD$ intersects $\overline{AC}$ at $E$ (distinct from $A$). The circumcircle of $\triangle ADC$ intersects $\overline{AB}$ at $F$ (distinct from $A$).

\begin{enumerate}[label=(\alph*),leftmargin=1.5cm]
  \item Prove that $BDEF$ is a cyclic quadrilateral.
  \item Prove that $AE = AF$.
  \item If $AB = AC = 10$ and $BC = 12$, find the radius of the circumcircle of $\triangle ABD$ when $D$ is the midpoint of $\overline{BC}$.
  \item Show that $\overline{EF} \parallel \overline{BC}$.
\end{enumerate}

\vfill
\begin{center}
\textit{--- End of Examination ---}
\end{center}

\newpage

% ============================================================
% ANSWER KEY
% ============================================================
\section*{Answer Key (For Teacher Use)}

\noindent\textbf{Section I: Multiple Choice (2 pts each)}
\begin{multicols}{2}
\bigskip
1. \textbf{B} --- Converse switches hypothesis and conclusion.\\[4pt]
2. \textbf{C} --- $180 - 52 - 68 = 60\deg$.\\[4pt]
3. \textbf{C} --- $1 + 2 = 3$, triangle inequality fails ($a+b > c$ required).\\[4pt]
4. \textbf{C} --- Contrapositive negates and reverses: ``If not a rhombus, then not a square.''\\[4pt]
5. \textbf{D} --- Dilation with $k \neq 1$ changes size; not a rigid motion.\\[4pt]
6. \textbf{B} --- Consecutive angles supplementary: $3x+15 + 2x+5 = 180 \Rightarrow 5x = 160 \Rightarrow x = 28$.\\[4pt]
7. \textbf{A} --- Area ratio $36\!:\!100 = 9\!:\!25$; linear ratio $3\!:\!5$; base $= 6 \cdot \frac{5}{3} = 10$.\\[4pt]
8. \textbf{B} --- $AB = \sqrt{9^2+12^2} = 15$; $\sin A = \frac{12}{15} = \frac{4}{5}$.\\[4pt]
9. \textbf{C} --- Central angle measure equals intercepted arc measure.\\[4pt]
10. \textbf{D} --- $2\pi r = 16\pi \Rightarrow r = 8$; $A = \pi(8)^2 = 64\pi$.\\[4pt]
11. \textbf{C} --- Circumcenter is intersection of perpendicular bisectors.\\[4pt]
12. \textbf{D} --- $V = \frac{4}{3}\pi(6)^3 = 288\pi$.\\[4pt]
13. \textbf{B} --- $5x-10 = 3x+20 \Rightarrow 2x = 30 \Rightarrow x = 15$.\\[4pt]
14. \textbf{A} --- $(x-h)^2 + (y-k)^2 = r^2 \Rightarrow (x+3)^2 + (y-4)^2 = 49$.\\[4pt]
15. \textbf{A} --- Perimeter ratio equals linear (similarity) ratio $= 3:5$.
\end{multicols}

\vspace{0.5cm}
\noindent\textbf{Section II: Short Answer (3 pts each)}
\begin{enumerate}[label=\arabic*.]
  \item $1$ (Pythagorean identity: $\cos^2\theta + \sin^2\theta = 1$)
  \item $AB = \sqrt{OA^2 - OB^2} = \sqrt{100 - 64} = \boxed{6~\text{cm}}$ (tangent $\perp$ radius)
  \item Apothem $= 6 \cdot \cot(30\deg) = \boxed{3\sqrt{3}}$
  \item Reflect $A(1,2)$ over $y$-axis $\to A'(-1,2)$; translate $3$ right $\to \boxed{A''(2,2)}$
  \item $ST = \frac{1}{2}QR = \boxed{7}$ (Midsegment Theorem)
  \item $\frac{1}{2}(230\deg) = \boxed{115\deg}$ (Inscribed angle theorem)
  \item $c = \sqrt{5^2+12^2} = 13$; altitude $= \frac{5 \cdot 12}{13} = \boxed{\frac{60}{13} \approx 4.62}$
  \item $SA = 2\pi r(r+h) = 2\pi(4)(4+10) = \boxed{112\pi}~\text{cm}^2$
  \item $m\angle P = \frac{1}{2}(110\deg - 50\deg) = \boxed{30\deg}$ (secant-secant angle)
  \item Lateral area $= 4 \cdot \frac{1}{2}(8)(7) = \boxed{112}$
\end{enumerate}

\vspace{0.3cm}
\noindent\textbf{Section III: Proofs (6 pts each)}
\begin{enumerate}[label=\arabic*.]
  \item[26.] Draw diagonal $\overline{AC}$. Show $\triangle ABC \cong \triangle CDA$ by SSS. Then $\angle BCA \cong \angle DAC$, so $AD \parallel BC$ (alt. int. $\angle$s). Similarly $AB \parallel DC$. Two pairs of parallel sides $\Rightarrow$ parallelogram.
  \item[27.] $\triangle ADB \cong \triangle EDC$ by SAS ($BD=DC$, $AD=DE$, vertical $\angle$s). So $AB=EC$. Similarly $AC=EB$. In $\triangle ABE$, $AB+AE > BE$. Substituting: $EC+AC > EB$. Since $AC > CD = \frac{1}{2}BC$, we get $BE+CE > BC$ by the triangle inequality.
  \item[28.] Suppose an acute angle $\angle A \geq 90\deg$. Then $\angle A + \angle C \geq 180\deg$, so $\angle B \leq 0\deg$, impossible. Thus $\angle A < 90\deg$.
  \item[29.] (1) Center compass at $A$, draw arc above and below $\overline{AB}$ with radius $> \frac{1}{2}AB$. (2) Same radius, center at $B$, draw arcs intersecting the first pair at points $P$ and $Q$. (3) Draw line $PQ$. Proof: $AP = BP = AQ = BQ$ (same radius), so $P$ and $Q$ are equidistant from $A$ and $B$, hence lie on the perpendicular bisector.
  \item[30.] Slopes: $m_{AB} = \frac{7-3}{4-(-2)} = \frac{4}{6} = \frac{2}{3}$; $m_{CD} = \frac{-3-1}{2-8} = \frac{-4}{-6} = \frac{2}{3}$. So $AB \parallel CD$. Similarly $m_{BC} = m_{DA} = -\frac{2}{3}$. Parallelogram confirmed. $AB = \sqrt{36+16} = \sqrt{52}$, $BC = \sqrt{16+64} = \sqrt{80}$. Adjacent sides not congruent $\Rightarrow$ not rhombus/square. $m_{AB} \cdot m_{BC} = -\frac{4}{9} \neq -1 \Rightarrow$ not rectangle. Conclusion: \textbf{parallelogram only}.
\end{enumerate}

\vspace{0.3cm}
\noindent\textbf{Section IV: Applications}
\begin{enumerate}[label=\arabic*.]
  \item[31.] $\sqrt{15^2 - 9^2} = \sqrt{144} = \boxed{12~\text{ft}}$
  \item[32.] $\frac{h}{24} = \frac{2}{3} \Rightarrow h = \boxed{16~\text{m}}$ (similar triangles)
  \item[33.] $R = 9$, $r = 6$; Area $= \pi(81 - 36) = \boxed{45\pi~\text{ft}^2}$
  \item[34.] $R^2 = 4^2 + (R-2)^2 \Rightarrow R^2 = 16 + R^2 - 4R + 4 \Rightarrow 4R = 20 \Rightarrow \boxed{R = 5~\text{ft}}$
  \item[35.] Centroid $= \frac{A+B+C}{3} \Rightarrow C = 3(4,5) - (1,2) - (7,8) = (12-8,\, 15-10) = \boxed{(4,\, 5)}$
  \item[36.] Area $= \frac{1}{2} \cdot 8 \cdot 10^2 \cdot \sin(45\deg) = 400 \cdot \frac{\sqrt{2}}{2} = \boxed{566}$
  \item[37.] $\frac{1}{3}\pi r^2(12) = 120\pi \Rightarrow r^2 = 30$; slant $l = \sqrt{30+144} = \boxed{\sqrt{174} \approx 13.2~\text{in}}$
  \item[38.] Altitude $= \sqrt{3 \cdot 12} = \boxed{6}$ (Geometric Mean Altitude Theorem)
\end{enumerate}

\vspace{0.3cm}
\noindent\textbf{Section V: Challenge Problems}
\begin{enumerate}
  \item[39.] (a) Central angle $= 360/12 = 30\deg$. By Law of Cosines: $s = \sqrt{R^2 + R^2 - 2R^2\cos 30\deg} = R\sqrt{2 - \sqrt{3}}$.
  \par(b) Area $= 12 \cdot \frac{1}{2} R^2 \sin 30\deg = 12 \cdot \frac{1}{2} R^2 \cdot \frac{1}{2} = \boxed{3R^2}$.
  \par(c) With $R = 10$: Area $= \boxed{300}$.
  \par(d) Area of inscribed $n$-gon $= \frac{1}{2} n R^2 \sin(2\pi/n)$. Let $x = 2\pi/n$. As $n \to \infty$, $x \to 0$. Area $= \pi R^2 \cdot \frac{\sin x}{x} \to \pi R^2$ since $\lim_{x\to 0} \frac{\sin x}{x} = 1$.
  \item[40.] (a) $\angle ADB + \angle ADC = 180\deg$ (linear pair). $\angle AEB = \angle ADB$ (same arc in first circle). $\angle AFC = \angle ADC$ (same arc in second circle). So $\angle AEB + \angle AFC = 180\deg$. Since $\angle BEF = 180\deg - \angle AEB$ and $\angle BFD = 180\deg - \angle AFC$, we get $\angle BEF + \angle BDF = 180\deg$, so $BDEF$ is cyclic.
  \par(b) Since $BDEF$ is cyclic, $\angle FEB = \angle FDB$. Since $\triangle ABC$ is isosceles, $\angle ABC = \angle ACB$. From the circumcircles, $\angle FDB = \angle FAB$ and $\angle EDC = \angle EAC$. By symmetry of the isosceles setup, $\triangle AFE \cong \triangle AEF'$, giving $AE = AF$.
  \par(c) $D$ midpoint $\Rightarrow BD = 6$, $AD = \sqrt{10^2 - 6^2} = 8$. In $\triangle ABD$: $AB=10$, $BD=6$, $AD=8$. Circumradius: $R = \frac{abc}{4K} = \frac{10 \cdot 6 \cdot 8}{4 \cdot 24} = \boxed{5}$.
  \par(d) Since $AE = AF$, $\triangle AEF$ is isosceles with $\angle AEF = \angle AFE = \frac{180 - \angle A}{2} = \angle B$. So $\angle AEF = \angle ABC$, giving $\overline{EF} \parallel \overline{BC}$ (corresponding angles).
\end{enumerate}

\end{document}
